% !TeX root = ../../Book4.tex
% 纲要标题：### 21.3 倾斜对象
% 纲要标签：b4:sec:2003
% 纲要位置：计划纲要.md，第 3039 行。

\section{倾斜对象}
\label{b4:sec:2003}

用复形代替正则模时，既要能够生成全部完美对象，又要消去非零移位的自身态射。三顶点例子的自同态计算将这两个条件落实，并产生新的导出等价。


% 纲要内容（仅作写作计划，尚非教材正文）：
%
% * 消失条件与生成条件
% * 经典有限维代数例子
% * 导出 Morita 理论的基本版本
%

\subsection{消失条件与生成条件}
\label{b4:subsec:2003:01}

第二卷定理16.3.3已经证明，普通 Morita 理论用有限生成投射生成元代替正则模。导出范畴允许把替代对象选成复形。假如希望一个三角等价把环 \(B\) 的正则模送到 \(R\) 上的完美复形 \(T\)，并在完美对象上给出等价，那么正则模的两项性质也须由 \(T\) 保留：\(B\) 与自身非零移位之间没有态射，而每个完美 \(B\)-复形都能由 \(B\) 经过有限次移位、锥和取直和项构造。因此需要
\[
 \operatorname{Hom}(T,T[n])=0\quad(n\ne0),\qquad
 \operatorname{thick}(T)=\operatorname{Perf}(R).
\]
前者使自同态信息集中在零次，后者保证用 \(T\) 能恢复全部完美对象。如果高次自身态射仍然存在，就须保留完整的 DG 自同态代数。

\ProseParagraphBreak
两个简单例子说明这两项要求各有什么作用。若 \(R=k\times k\)，直和项 \(T=k\times0\) 没有非零高次自身态射，但由它构造的对象都被 \((0,1)\) 零化，因而无法得到 \(0\times k\)。若 \(R=k\)，\(T=k\oplus k[1]\) 包含 \(k\) 作为直和项，所以生成全部完美复形；但从 \(T\) 的 \(k[1]\) 分量到 \(T[1]\) 的 \(k[1]\) 分量可取恒等映射，故 \(\operatorname{Hom}(T,T[1])\ne0\)。

\begin{definition}\label{b4:def:tilting-complex}
设 \(k\) 为域，\(R\) 为含幺结合 \(k\)-代数。若 \(T\in\operatorname{Perf}(R)\) 满足
\[
 \operatorname{Hom}_{D(R\text{-}\mathrm{Mod})}(T,T[n])=0
 \quad(n\ne0),\qquad
 \operatorname{thick}(T)=\operatorname{Perf}(R),
\]
称 \(T\) 为 \(R\) 上的\term[qingxiefuxing]{倾斜复形}[tilting complex]，也称所给三角范畴中的倾斜对象。这里使用的是包括正负所有非零次数的消失条件。
\end{definition}

\(R\) 本身是倾斜复形：其自身高次态射为零，厚生成条件由\cref{b4:thm:perfect-compact-thick}给出。倾斜复形推广的正是正则模的这两项性质。

\ProseParagraphBreak
这里的倾斜对象与\chapref{b4:ch:19}的扭对倾斜使用不同数据。后者从一个心及其中的扭对构造另一个心；本节则要求一个完美对象具有自身态射消失和厚生成性质。构造出新心并不自动产生这样的对象，更不自动给出两个模导出范畴之间的等价。

\begin{proposition}\label{b4:prop:tilting-module-criterion}
设 \(k\) 为域，\(R\) 为 \(k\)-代数，\(T\) 为左 \(R\)-模。假设 \(T\) 有有限长有限生成投射消解，\(\operatorname{Ext}^n_R(T,T)=0\) 对所有 \(n>0\) 成立，并且存在正合列
\[
 0\longrightarrow R\longrightarrow T_0\longrightarrow T_1
 \longrightarrow\cdots\longrightarrow T_m\longrightarrow0,
\]
其中每个 \(T_i\) 是某个有限直和 \(T^{r_i}\) 的直和项。则 \(T\) 集中在零次是倾斜复形。
\end{proposition}
\begin{proof}
有限投射消解给出完美性。对 \(n>0\)，自身导出态射由\cref{b4:thm:derived-ext-identification}识别为所给 Ext 群；对 \(n<0\)，取位于非正次数的投射消解 \(P\to T\)，则 \(\underline{\operatorname{Hom}}_R(P,T)\) 没有负次项，故负次态射为零。

把所给长正合列按像分成短正合列。从末项起逆向使用相应好三角，可知 \(R\in\operatorname{thick}(T)\)。由\cref{b4:thm:perfect-compact-thick}，\(\operatorname{Perf}(R)=\operatorname{thick}(R)\subseteq\operatorname{thick}(T)\)。反向包含由完美子范畴的厚性和 \(T\) 完美得到。
\end{proof}

\subsection{经典有限维代数例子}
\label{b4:subsec:2003:02}

下面用三个向量空间及两个线性映射计算一个倾斜对象，暂不需要第五卷的一般箭图表示论。设 \(k\) 为域，\(R\) 是箭图 \(1\to2\to3\) 的路径代数：以三条长度零的路径、两条箭头及它们的复合为 \(k\)-基，乘法按可接续路径复合，不可接续时取零。具体记箭头为 \(\alpha:1\to2\)、\(\beta:2\to3\)，约定 \(\alpha=e_2\alpha e_1\)、\(\beta=e_3\beta e_2\)，所以 \(\beta\alpha\) 是从一到三的路径。有限维左 \(R\)-模等价于图式 \(V_1\to V_2\to V_3\)；三个顶点幂等元给出三个分量，箭头作用给出两个线性映射。模同态是使这两个方块交换的三元线性映射。

\ProseParagraphBreak
写出三个投射模以及中间顶点的简单模：
\[
\begin{aligned}
 P_1&=(k\xrightarrow{1}k\xrightarrow{1}k),&
 P_2&=(0\longrightarrow k\xrightarrow{1}k),\\
 P_3&=(0\longrightarrow0\longrightarrow k),&
 S_2&=(0\longrightarrow k\longrightarrow0).
\end{aligned}
\]
这里 \(P_i=Re_i\)，并且 \(\operatorname{Hom}_R(P_i,V)\cong V_i\)：态射由顶点 \(i\) 处的向量决定，其余分量由箭头推出。因此它们确为投射模，且 \(R=P_1\oplus P_2\oplus P_3\)。

\begin{example}\label{b4:exm:a3-tilting-object}
在上述路径代数 \(R\) 上，令 \(T=P_1\oplus P_2\oplus S_2\)。则 \(T\) 是倾斜复形，而
\[
 B=\operatorname{End}_R(T)^{\mathrm{op}}
\]
同构于路径代数 \(k(1\to2\leftarrow3)\)。
\end{example}
\begin{proof}
有短正合列
\begin{equation}\label{b4:eq:a3-simple-resolution}
 0\longrightarrow P_3\xrightarrow{i}P_2
 \xrightarrow{q}S_2\longrightarrow0.
\end{equation}
所以 \(T\) 有长度至多一的有限投射消解。对任意表示 \(V\)，这份消解给出
\[
 \operatorname{Ext}^1_R(S_2,V)
 =\operatorname{coker}(V_2\longrightarrow V_3),
 \qquad \operatorname{Ext}^{n}_R(S_2,V)=0\quad(n\ge2).
\]
对 \(V=P_1,P_2,S_2\)，\(V_2\to V_3\) 都满；因此 \(\operatorname{Ext}^n_R(T,T)=0\) 对所有 \(n>0\) 成立。负次态射为零同样由非正次投射消解计算。\eqref{b4:eq:a3-simple-resolution}的三角把 \(P_3\) 从 \(P_2,S_2\) 构造出来，故 \(R\in\operatorname{thick}(T)\)，由完美性即得厚生成。

自同态的计算也可以逐项完成。三个直和项的自同态环均为 \(k\)。不同项之间仅有
\[
 \operatorname{Hom}_R(P_2,P_1)=k j,
 \qquad \operatorname{Hom}_R(P_2,S_2)=k q
\]
非零：从 \(P_1\) 出发的态射由目标顶点一的分量决定，故另外两项为零；从 \(S_2\) 到 \(P_1\) 或 \(P_2\) 的态射则由第二条箭头的交换条件迫为零。相对于 \(T\) 的三个直和项，自同态代数因此为
\[
 E=\left\{\left.
 \begin{pmatrix}a&b&0\\0&c&0\\0&d&e\end{pmatrix}
 \ \right|\ a,b,c,d,e\in k\right\}.
\]
其中两条非恒等箭头从第二个直和项发出，彼此不能复合。取反代数后，两条箭头都指向第二个顶点，正好得到 \(B=k(1\to2\leftarrow3)\)。
\end{proof}

这个例子不是普通 Morita 等价换了一种写法。\(R\) 的三个不可分解投射模的合成长度为 \(3,2,1\)；\(B\) 的三个不可分解投射模分别支持在 \(1\to2\)、顶点二、\(3\to2\)，合成长度为 \(2,1,2\)。这些类型已经穷尽：两代数模去箭头生成的幂零理想均为 \(k^3\)，三个顶点给出全部简单模，第二卷定理15.3.10的简单模与不可分解投射模对应给出所述清单。模范畴等价保持简单对象、投射性、不可分解性和合成列，因而保持这组长度；两组数不同，所以两者的模范畴不等价。下一小节的定理却会给出其导出范畴的等价。

\subsection{导出 Morita 理论的基本版本}
\label{b4:subsec:2003:03}

令 \(T\) 为倾斜复形，选择有界有限生成投射模型 \(P\)，并令
\[
 E=\underline{\operatorname{End}}_R(P),\qquad
 B=\operatorname{End}_{D(R\text{-}\mathrm{Mod})}(T)^{\mathrm{op}}.
\]
倾斜条件恰好说 \(H^n(E)=0\) 对 \(n\ne0\) 成立，而 \(H^0(E)=B^{\mathrm{op}}\)。这一步把一个涉及所有移位的条件变为自同态 DG 代数的上同调消失。

\begin{proposition}\label{b4:prop:dg-degree-zero-formality}
设 \(k\) 为交换环，\(E\) 为 \(k\) 上的 DG 代数，且 \(H^n(E)=0\) 对所有 \(n\ne0\) 成立。令
\[
 (\tau_{\le0}E)^n=
 \begin{cases}
 E^n,&n<0,\\
 \ker(\mathrm{d}:E^0\to E^1),&n=0,\\
 0,&n>0.
 \end{cases}
\]
原微分和乘法使 \(\tau_{\le0}E\) 成为 DG 子代数，而且包含映射与零次取商给出 DG 代数拟同构
\[
 E\ \longleftarrow\ \tau_{\le0}E\ \longrightarrow\ H^0(E).
\]
\end{proposition}
\begin{proof}
非正次数在乘法下封闭；两零次余圈之积仍为余圈。来自 \(E^{-1}\) 的微分落在零次余圈中，故这是子复形。到 \(H^0(E)\) 的映射在负次取零，在零次把余圈取模余边。负次元素的乘积不会产生正次或零次项；两个零次余圈的乘法与取商相容，所以它是代数映射。包含映射在非正次上同调上为同构，而 \(E\) 正次上同调为零；取商映射在零次为同构，两边其他次数上同调均为零。因此两映射都是拟同构。
\end{proof}

注意这里得到的是通过第三个 DG 代数连接的两个拟同构，未必有直接代数映射 \(H^0(E)\to E\)。任选每个同伦类的复形映射代表，通常不能同时保持乘法。正是这一点阻止了把普通 Morita 理论中的论证逐字照搬。

\begin{theorem}[导出 Morita 定理的倾斜版本]\label{b4:thm:derived-morita-tilting}
设 \(k\) 为域，\(R\) 为含幺结合 \(k\)-代数，\(T\in\operatorname{Perf}(R)\) 为倾斜复形，并令
\[
 B=\operatorname{End}_{D(R\text{-}\mathrm{Mod})}(T)^{\mathrm{op}}.
\]
则存在 \(R\)-\(B\) 双模复形 \(X\)，作为左 \(R\)-模复形时在导出范畴中同构于 \(T\)，使
\[
 X\otimes_B^{\mathbf L}-:
 D(B\text{-}\mathrm{Mod})\ \longrightarrow\ D(R\text{-}\mathrm{Mod})
\]
成为 \(k\)-线性三角等价，且把 \(B\) 送到 \(T\)。此等价限制为 \(\operatorname{Perf}(B)\simeq\operatorname{Perf}(R)\)。
\end{theorem}

这一定理是\term[daochuMoritalilun]{导出 Morita 定理}[derived Morita theorem]的倾斜版本。下面通过 DG 自同态代数说明双模及等价的构造；所用文献主要采用右模，本节使用左模，因此自同态代数取反。

\begin{proofsketch}
取 \(T\) 的有界有限生成投射模型 \(P\)，记
\(E=\underline{\operatorname{End}}_R(P)\)、\(C=E^{\mathrm{op}}\)。这里取 DG 反代数：对次数分别为 \(a,b\) 的齐次自同态 \(e,f\)，
\[
e^{\mathrm{op}}f^{\mathrm{op}}=(-1)^{ab}(fe)^{\mathrm{op}},
\qquad \mathrm d_C(e^{\mathrm{op}})=(\mathrm d_Ee)^{\mathrm{op}}.
\]
在 \(P\) 上定义右作用
\(p\cdot e^{\mathrm{op}}=(-1)^{a|p|}e(p)\)。反乘法中的符号保证右作用的结合律；Hom 微分公式则给出
\[
\mathrm d_P(p\cdot e^{\mathrm{op}})
=(\mathrm d_Pp)\cdot e^{\mathrm{op}}
+(-1)^{|p|}p\cdot\mathrm d_C(e^{\mathrm{op}}).
\]
左 \(R\) 作用与这项右作用交换，所以 \(P\) 是严格的左 \(R\)、右 DG \(C\) 双模。

左 DG \(C\) 模是带分次左 \(C\) 作用的复形 \(N\)，其微分满足 \(\mathrm d(cn)=(\mathrm d_Cc)n+(-1)^{|c|}c\,\mathrm dn\)。以下 \(D(C)=D(C\text{-}\mathrm{DGMod})\) 表示左 DG \(C\) 模同伦范畴在拟同构处局部化所得的导出范畴。对左 \(R\)-模复形 \(Y\)，次数 \(b\) 的齐次映射 \(\varphi:P\to Y\) 上的左 \(C\) 作用为
\[
e^{\mathrm{op}}\cdot\varphi=(-1)^{ab}\varphi e.
\]
带符号的反乘法给出作用的结合律，而
\(\mathrm d(\varphi e)=(\mathrm d\varphi)e+(-1)^b\varphi(\mathrm d_Ee)\) 给出
\[
\mathrm d(e^{\mathrm{op}}\cdot\varphi)
=\mathrm d_C(e^{\mathrm{op}})\cdot\varphi
+(-1)^a e^{\mathrm{op}}\cdot\mathrm d\varphi.
\]
因此 Hom 复形确实是左 DG \(C\) 模。对这些 DG 模使用半自由消解，得到伴随
\[
 F=P\otimes_C^{\mathbf L}-:D(C)\rightleftarrows
 D(R\text{-}\mathrm{Mod}):U=\mathbf R\operatorname{Hom}_R(P,-).
\]
因为 \(P\) 有界且逐项有限生成投射，\(U\) 可用普通 Hom 复形计算，并保持任意余积；左伴随 \(F\) 也保持任意余积。复形级伴随的单位在齐次元上为
\(\eta_N(n)(p)=(-1)^{|n||p|}p\otimes n\)。在识别 \(P\otimes_CC\cong P\) 后，它在左正则 DG 模 \(C\) 上的公式是
\[
\eta_C:C\longrightarrow\underline{\operatorname{End}}_R(P),
\qquad
\eta_C(e^{\mathrm{op}})(p)
=(-1)^{a|p|}p\cdot e^{\mathrm{op}}=e(p).
\]
它与微分相容，并且
\[
\eta_C(e^{\mathrm{op}}f^{\mathrm{op}})
=(-1)^{ab}fe
=e^{\mathrm{op}}\cdot\eta_C(f^{\mathrm{op}}).
\]
故这是左 \(C\) 模复形的同构，目标使用刚才的左作用。它并不是把反代数与原自同态代数当作同一 DG 代数。

单位为同构的对象组成一个在同构下封闭的满子范畴。两函子均为三角函子，单位与平移相容，三角的 Hom 正合列给出该类对平移及好三角封闭；两函子保持任意余积，故该类也对任意余积封闭；加性与自然性又使单位在直和分解下分块，故该类对直和项封闭。它包含 \(C\)，而半自由消解把每个 DG 模写成自由移位的逐层锥及望远镜，所以它等于整个 \(D(C)\)。这证明 \(F\) 全忠实。

由伴随三角恒等式，余单位在 \(F(C)\cong P\) 上也是同构。余单位为同构的对象类同样在同构下封闭，并由上述三角、余积和分块论证分别对好三角、任意余积和直和项封闭。倾斜条件使
\(R\in\operatorname{thick}(P)\)，故该类包含 \(R\)。以
\(\operatorname{Loc}(R)\) 表示包含 \(R\)、对平移、好三角、任意余积及直和项封闭的最小满子范畴。普通环的半自由消解与望远镜构造给出
\[
\operatorname{Loc}(R)=D(R\text{-}\mathrm{Mod}).
\]
因此余单位处处为同构，\(F\) 是等价。这里先用有限构造得到 \(R\)，再用任意余积与望远镜恢复无界对象。

自身态射消失使 \(H^{\ne0}(C)=0\)、\(H^0(C)=B\)。由\cref{b4:prop:dg-degree-zero-formality}，令 \(A=\tau_{\le0}C\)，便有 DG 代数拟同构
\[
 C\longleftarrow A\longrightarrow B.
\]
每条拟同构都诱导导出模范畴的等价：在自由模上，延标量与限制标量的伴随单位是这条拟同构；再用半自由消解及余积、三角封闭性推广到全部对象。限制标量反映拟同构，伴随三角恒等式便使余单位也处处为同构。合成上述等价，便得到 \(D(B)\simeq D(R)\)。

严格双模也由这一步得到。把 \(P\) 限制为 \(R\)-\(A\) 双模，取它在 DG 代数 \(R\otimes_kA^{\mathrm{op}}\) 上的半自由替换 \(Q\to P\)，并令
\[
 X=Q\otimes_A B.
\]
设 \(\mathcal B_R\) 是 \(R\) 的 \(k\)-向量空间基。\(Q\) 的半自由滤过逐层在分次模意义下分裂，各滤过商是移位正则模
\(R\otimes_kA^{\mathrm{op}}\) 的直和。由于 \(R\) 集中在零次，限制为左 \(A^{\mathrm{op}}\) 模后有复形同构
\[
R\otimes_kA^{\mathrm{op}}
\cong\bigoplus_{r\in\mathcal B_R}A^{\mathrm{op}}.
\]
因此限制后的滤过商仍是移位自由 DG \(A^{\mathrm{op}}\) 模的直和；换成右 \(A\) 模即说明 \(Q\) 仍半自由，因而 K-平坦。将 \(Q\otimes_A-\) 作用于拟同构 \(A\to B\)，得到左 \(R\) 模复形的拟同构
\(Q\cong Q\otimes_AA\to Q\otimes_AB=X\)。于是
\(X\simeq P\simeq T\) 作为左 \(R\) 模成立。导出张量的结合律把合成等价识别为 \(X\otimes_B^{\mathbf L}-\)。它把 \(B\) 送到 \(T\)，从而把 \(\operatorname{thick}(B)\) 送到 \(\operatorname{thick}(T)\)；\cref{b4:thm:perfect-compact-thick}给出完美子范畴上的限制。

本概要省略了 DG 半自由消解的逐次添加生成元、双模替换的比较及导出伴随的链级相容性。消解与双模张量分别见\cite[Sections 3.1 and 6.1--6.3]{Keller1994DerivingDG}，紧生成下的等价判据及倾斜实现见\cite[Sections 4.2, 8.2 and 9.1--9.2]{Keller1994DerivingDG}；普通代数的同一定理见\cite[Section 4, Theorem (Rickard)]{Keller1998AbelianDerived}。这些构造把同伦意义的自同态作用换成严格双模作用，不能用逐项任取代表代替。
\end{proofsketch}

把定理用于\cref{b4:exm:a3-tilting-object}，便得到 \(k(1\to2\to3)\) 与 \(k(1\to2\leftarrow3)\) 的导出等价。虽然它们的模范畴不同，允许使用复形后，原来的投射生成元可以由 \(P_1\oplus P_2\oplus S_2\) 代替；\eqref{b4:eq:a3-simple-resolution}正是恢复缺少的投射模 \(P_3\) 的计算。
